\subsection{诱导表示与余诱导表示}

设 H 是群 G 的一个子群。

若 \((V,\pi)\) 是 G 的一个线性表示，则 \(\pi\) 在 H 上的限制是 H 在 V 中的一个线性表示；我们把它记作 \(\operatorname{Res}_{\mathrm{H}}^{\mathrm{G}}(\pi)\)。K{[}H{]}-模（与 \(\operatorname{Res}_{\mathrm{H}}^{\mathrm{G}}(\pi)\) 相伴的那个模）就是由 K{[}G{]}-模 V 经标量限制得到的模（II, §1, No.~13, 第 221 页）。若 V 是有限维自由 K-模，则 \(\operatorname{Res}_{\mathrm{H}}^{\mathrm{G}}(\pi)\) 的特征标是 \(\pi\) 的特征标在 H 上的限制。

设 \((\mathrm{M},\sigma)\) 是 H 的一个表示。

我们把 K{[}G{]} 视为一个 (K{[}G{]}, K{[}H{]})-双模，把 M 视为一个 K{[}H{]}-模。K{[}G{]}-模 \(\mathcal{T}(M) = K[G] \otimes_{K[H]} M\)（VIII, 第 58 页）定义了 G 的一个线性表示，记作 \(\operatorname{Ind}_{\mathrm{H}}^{\mathrm{G}}(\sigma)\)，称为由 \(\sigma\) 诱导的 G 的表示。若 \((\mathrm{V}, \pi)\) 是 G 的一个线性表示，则 K{[}H{]}-模 \(\mathcal{H}(\mathrm{V}) = \mathrm{Hom}_{\mathrm{K[G]}}(\mathrm{K[G]},\mathrm{V})\) 可与一个 K{[}H{]}-模等同，该模对应于表示 \(\mathrm{Res}_{\mathrm{H}}^{\mathrm{G}}(\pi)\)。于是，伴随态射（VIII, 第 59 页）给出一个从 \(\mathrm{Hom}_{\mathrm{H}}(\sigma ,\mathrm{Res}_{\mathrm{H}}^{\mathrm{G}}(\pi))\) 到 \(\mathrm{Hom}_{\mathrm{G}}(\mathrm{Ind}_{\mathrm{H}}^{\mathrm{G}}(\sigma),\pi)\) 的 K-模同构，称为典范同构（"弗罗贝尼乌斯互反律"）。

我们把 K{[}G{]} 视为一个 (K{[}H{]}, K{[}G{]})-双模。K{[}G{]}-模 \(\mathcal{H}(M) = \operatorname{Hom}_{\mathrm{K}[\mathrm{H}]}(\mathrm{K}[\mathrm{G}], \mathrm{M})\) 定义了 G 的一个表示，记作 \(\operatorname{Coind}_{\mathrm{H}}^{\mathrm{G}}(\sigma)\)，称为由 \(\sigma\) 余诱导的 G 的表示。若 \((\mathrm{V}, \pi)\) 是 G 的一个线性表示，则 K{[}H{]}-模 \(\mathcal{T}(\mathrm{V}) = \mathrm{K}[\mathrm{G}] \otimes_{\mathrm{K}[\mathrm{G}]} \mathrm{V}\) 可与一个 K{[}H{]}-模等同，该模对应于表示 \(\operatorname{Res}_{\mathrm{H}}^{\mathrm{G}}(\pi)\)。于是，伴随态射（同前引文）给出一个从 \(\operatorname{Hom}_{\mathrm{H}}(\operatorname{Res}_{\mathrm{H}}^{\mathrm{G}}(\pi), \sigma)\) 到 \(\operatorname{Hom}_{\mathrm{G}}(\pi, \operatorname{Coind}_{\mathrm{H}}^{\mathrm{G}}(\sigma))\) 的 K-模同构，称为典范同构。

设 \(\varepsilon : \mathrm{K}[\mathrm{G}] \to \mathrm{K}[\mathrm{H}]\) 是由关系式 \(\varepsilon(h) = h\)（若 \(h \in \mathrm{H}\)）与 \(\varepsilon(g) = 0\)（若 \(g \in \mathrm{G} - \mathrm{H}\)）刻画的 K-模同态。映射 \(\varepsilon\) 是一个 \((\mathrm{K}[\mathrm{H}], \mathrm{K}[\mathrm{H}])\)-双模同态。设 \((\mathrm{M}, \sigma)\) 是 H 的一个线性表示。映射 \(v \mapsto v \circ \varepsilon\)（从 \(\mathrm{Hom}_{\mathrm{K}[\mathrm{H}]}(\mathrm{K}[\mathrm{H}], \mathrm{M})\) 到 \(\mathrm{Hom}_{\mathrm{K}[\mathrm{H}]}(\mathrm{K}[\mathrm{G}], \mathrm{M})\)）是一个 \(\mathrm{K}[\mathrm{H}]\)-模同态。把 M 与 \(\mathrm{Hom}_{\mathrm{K}[\mathrm{H}]}(\mathrm{K}[\mathrm{H}], \mathrm{M})\) 等同，我们就得到一个 \(\mathrm{K}[\mathrm{H}]\)-模同态，从 M 到 \(\mathrm{Res}_{\mathrm{H}}^{\mathrm{G}}(\mathrm{Coind}_{\mathrm{H}}^{\mathrm{G}}(\sigma))\)。弗罗贝尼乌斯互反律把它变为一个模同态 \(\iota\)（\(\mathrm{K}[\mathrm{G}]\)-模的同态），从 \(\mathrm{Ind}_{\mathrm{H}}^{\mathrm{G}}(\sigma)\) 到 \(\mathrm{Coind}_{\mathrm{H}}^{\mathrm{G}}(\sigma)\)，它由关系式

\[
\iota (g \otimes m) (g ^ {\prime}) = \varepsilon (g ^ {\prime} g) m
\]

（对 \(g, g' \in G\) 与 \(m \in M\)）刻画。这个同态称为典范同态。

命题 2. —— 设 H 是 G 的一个子群，\((\mathrm{M},\sigma)\) 是 H 的一个线性表示。典范同态 \(\iota:\operatorname{Ind}_{\mathrm{H}}^{\mathrm{G}}(\sigma)\to\operatorname{Coind}_{\mathrm{H}}^{\mathrm{G}}(\sigma)\) 是单射。若子群 H 在 G 中有有限指数，则 \(\iota\) 是一个 K{[}G{]}-模同构。

设 \(S \subset G\) 是 \(G/H\) 的一个代表系。族 \((s)_{s \in S}\) 是右 \(K[H]\)-模 \(K[G]\) 的一组基。由此得出，映射 \(M^{(S)} \to \operatorname{Ind}_{H}^{G}(\sigma)\)（由 \((m_s)_{s \in S} \mapsto \sum_{s \in S} s \otimes m_s\) 定义）是一个 \(K\)-模同构。对一切 \(s, s' \in S\) 与每个 \(m \in M\)，我们有关系式

\[
\iota (s ^ {\prime} \otimes m) (s ^ {- 1}) = \left\{ \begin{array}{l} m \text {   if   } s = s ^ {\prime}, \\ 0 \text {   otherwise }. \end{array} \right.
\]

由此得出 \(\iota'\) 是单射。

设 H 在 G 中有有限指数。此时集合 S 是有限的；设 \(\rho\) 是从 \(\operatorname{Coind}_{\mathrm{H}}^{\mathrm{G}}(\sigma)\) 到 \(\operatorname{Ind}_{\mathrm{H}}^{\mathrm{G}}(\sigma)\) 的映射，由 \(u \mapsto \sum_{s \in S} s \otimes u(s^{-1})\) 给出。它满足 \((\iota \circ \rho(u))(s^{-1}) = u(s^{-1})\)（对 \(u \in \operatorname{Coind}_{\mathrm{H}}^{\mathrm{G}}(\sigma)\) 与 \(s \in S\)）。由于族 \((s^{-1})_{s\in\mathrm{S}}\) 是左 K{[}H{]}-模 K{[}G{]} 的一组基，映射 \(\iota\circ\rho\) 是恒同映射。于是 \(\iota\) 是双射，其逆双射为 \(\rho\)。

设 H 是 G 的一个有限指数子群。设 u 是 H 上的一个中心函数；用 \(u^{0}\) 表示 G 上延拓 u 且在 G − H 的每一点取零的函数。设 S 是 G/H 的一个代表系。对 \(g \in G\)，令

\[
\operatorname{Ind} _ {\mathrm{H}} ^ {\mathrm{G}} (u) (g) = \sum_ {s \in \mathrm{S}} u ^ {0} (s ^ {- 1} g s).\tag{4}
\]

注意，对一切 \(x, g \in G\) 与每个 \(h \in H\)，有 \(u^{0}((xh)^{-1}gxh) = u^{0}(x^{-1}gx)\)。由此得出，\(\operatorname{Ind}_{\mathrm{H}}^{\mathrm{G}}(u)\) 是 G 上的一个不依赖于 S 的选取的中心函数。这样我们就定义了一个 K-线性映射 \(Ind_{H}^{G}\)，从 \(\mathcal{Z}_{\mathrm{K}}(\mathrm{H})\) 到 \(\mathcal{Z}_{\mathrm{K}}(\mathrm{G})\)。

命题 3. —— 设 H 是 G 的一个有限指数子群。设 \((M,\sigma)\) 是 H 的一个线性表示。假设 M 是一个有限维自由 K-模。用 \((V,\pi)\) 表示由 \(\sigma\) 诱导的 G 的表示。那么 K-模 V 是自由且有限维的，并且

\[
\chi_ {\pi} = \mathrm{Ind} _ {\mathrm{H}} ^ {\mathrm{G}} (\chi_ {\sigma}).\tag{5}
\]

设 S 是 G/H 的一个代表系。如上所见，线性映射 \(\mathrm{M}^{\mathrm{S}}\to\mathrm{Ind}_{\mathrm{H}}^{\mathrm{G}}(\mathrm{M})\)（由 \((m_{s})_{s\in\mathrm{S}}\mapsto\sum_{s\in\mathrm{S}}s\otimes m_{s}\) 给出）是一个 K-模同构。特别地，V 是一个有限维自由 K-模。对任意 \(g\in G\)，用 \(M_{g}\) 表示 M 在映射 \(m\mapsto g\otimes m\) 下的像。对一切 \(g,g'\in G\)，\(M_{g}=M_{g'}\) 当且仅当 \(gH=g'H\)，且 V 是诸子模 \(M_{s}\)（\(s\in S\)）的直和。对一切 \(g,g'\in G\)，我们还有 \(\pi(g)\mathrm{M}_{g'}=\mathrm{M}_{gg'}\)。对任意 \(g\in G\)，用 \(S_{g}\) 表示 S 中满足 \(s^{-1}gs\in H\) 的元素 s 所成的集合。对一切 \(g,g'\in G\)（满足 \(g'^{-1}gg'\in H\)），V 的自同构 \(\pi(g)\) 诱导一个自同构 \(\pi(g)_{g'}\)（\(M_{g'}\) 的自同构），并且我们有

\[
\operatorname{Tr} (\pi (g)) = \sum_ {s \in \mathrm{S} _ {g}} \operatorname{Tr} (\pi (g) _ {s}) = \sum_ {s \in \mathrm{S} _ {g}} \operatorname{Tr} (\sigma (s ^ {- 1} g s)).
\]

命题的最后一个断言由此得出。

